\chapter{روش انجام پروژه}
\section{مقدمه}
در این فصل روش انجام پروژه به طور کامل شرح داده میشود. حداکثر صفحات این فصل 20 صفحه است.

\section{محتوا}
جدول (2-1) به صورت نمونه ارائه شده است.

\begin{table}[h!]
    \centering
    \caption{نتيجه بررسي پرسش نامه ها در ارتباط با عوامل موثر \cite{sample_ref}}
    \label{tab:factors}
    \renewcommand{\arraystretch}{1.2}
    \begin{tabular}{c l c}
        \toprule
        \textbf{رديف} & \textbf{عوامل موثر}     & \textbf{درصد} \\
        \midrule
        1             & احساس تعلق به سازمان    & 1/95          \\
        2             & نقش مديريت سازمان       & 7/87          \\
        3             & عوامل درون سازماني      & 9/82          \\
        4             & برگزاري دوره هاي آموزشي & 9/82          \\
        \bottomrule
    \end{tabular}
\end{table}

شکل (2-1) نیز به صورت نمونه ارائه شده است.

\begin{figure}[h!]
    \centering
    \includegraphics[width=0.6\textwidth]{example-image}
    \caption{نمونه شکل}
    \label{fig:sample}
\end{figure}

\noindent
متن
\\[0.3cm]
كه در آن $X$ $(\text{kg/m}^3)$ غلظت توده سلولي، $\mu$ $(\text{h}^{-1})$ شدت رشد ويژه و $D$ $(\text{h}^{-1})$ شدت رقيق‌سازي ميباشد.

\subsection{علت انتخاب روش}
دليل يا دلايل انتخاب روش انجام پروژه را تشريح مي‌کند.
در زير به تعدادي از روش‌هاي انجام پروژه اشاره شده است:
\begin{itemize}
    \item \textbf{روش آزمايشگاهي‌:} توصيف كامل برنامه‌ي آزمايشگاهي شامل مواد مصرفي و نحوه‌ي ساخت نمونه‌ها، شرح آزمايش‌ها شامل نحوه تنظيم و آماده سازي آزمايش‌ها و دستگاه‌هاي مورد استفاده و دقت و نحوه‌ي كاليبره كردن، شرح دستگاه ساخته شده (در صورت ساخت) و ارائه‌ي روش اعتبارسنجي.
    \item \textbf{روش آماري:} توصيف ابزارهاي گردآوري اطلاعات كمي و كيفي، اندازه‌ي نمونه‌ها، روش نمونه‌برداري، تشريح مباني روش آمار و ارائه‌ي روش اعتبارسنجي.
    \item \textbf{روش نرم‌افزارنويسي:} توصيف كامل برنامه‌نويسي، مباني برنامه و ارائه‌ي روش اعتبارسنجي.
    \item \textbf{روش مطالعه‌ي موردي:} توصيف كامل محل و موضوع مطالعه، علت انتخاب مورد و پارامترهايي كه تحت ارزيابي قرار داده مي‌شوند، و ارائه‌ي روش اعتبارسنجي.
    \item \textbf{روش تحليلي يا مدلسازي:} توصيف كامل مباني يا اصول تحليل يا مدل و ارائه‌ي روش اعتبارسنجي.
    \item \textbf{روش ميداني:} چگونگي دستيابي به دادهها در ميدان عمل و نحوه برداشت از پاسخهاي دريافتي.
\end{itemize}
\subsection{تشريح كامل روش انجام پروژه}
% +++++++++++++++++++++++++++++++++++++++++++++++++++++++++++++++++++++++++++++++++++++++++++++++++
\section{معرفی سیستم}

در این سیستم توسعه‌یافته، یک آونگ ساده به ارابه‌ای متصل است که روی یک ریل حرکت می‌کند. طول کابل آونگ متغیر بوده و به عنوان یک متغیر سینماتیکی (ورودی) تحت عنوان $l(t)$ در نظر گرفته می‌شود. همچنین سیستم توسط دو نیروی محرک $F_1$ و $F_2$ از دو انتهای ریل تحریک می‌گردد (هر دو در جهت استاندارد محور $x$ مثبت فرض شده‌اند).

مختصات تعمیم‌یافته سیستم با توجه به دو درجه آزادی موجود به صورت زیر انتخاب می‌شوند:
\[
    q = \begin{bmatrix} x \\ \theta \end{bmatrix}
\]
که در آن:
\begin{itemize}
    \item $x$ : مکان افقی ارابه نسبت به تکیه‌گاه اول
    \item $\theta$ : زاویه آونگ نسبت به راستای قائم (پادساعت‌گرد مثبت است)
    \item $l(t)$ : طول متغیر کابل آونگ (متغیر ورودی سینماتیکی)
    \item $M$ : جرم ارابه
    \item $m$ : جرم نقطه‌ای آونگ
    \item $M_b$ : جرم ریل با طول کل $D$
    \item $g$ : شتاب گرانش
\end{itemize}

اصطکاک‌های لزج سیستم:
\begin{itemize}
    \item $b_c$ : ضریب اصطکاک لزج ارابه و ریل
    \item $b_p$ : ضریب اصطکاک لزج مفصل آونگ
\end{itemize}

خروجی‌های مطلوب شامل موقعیت‌ها و نیروهای تکیه‌گاهی ریل ($R_{x1}, R_{y1}$ در تکیه‌گاه اول و $R_{y2}$ در تکیه‌گاه دوم) می‌باشند.

\section{استخراج معادلات حرکت به روش لاگرانژ}

\subsection{سینماتیک و مختصات جرم آونگ}

با رعایت قانون دست راست، محور $x$ به سمت راست و محور $y$ به سمت بالا در نظر گرفته شده است. مختصات جرم آونگ برابر است با:
\[
    x_p = x + l(t) \sin\theta
\]
\[
    y_p = -l(t) \cos\theta
\]
با مشتق‌گیری نسبت به زمان ($\dot{l} \neq 0$):
\[
    \dot{x}_p = \dot{x} + \dot{l}\sin\theta + l\dot{\theta}\cos\theta
\]
\[
    \dot{y}_p = -\dot{l}\cos\theta + l\dot{\theta}\sin\theta
\]
مربع سرعت خطی جرم آونگ ($v_p^2 = \dot{x}_p^2 + \dot{y}_p^2$) پس از بسط و ساده‌سازی روابط برابر است با:
\[
    v_p^2 = \dot{x}^2 + \dot{l}^2 + l^2\dot{\theta}^2 + 2\dot{x}\dot{l}\sin\theta + 2\dot{x}l\dot{\theta}\cos\theta
\]

\subsection{انرژی جنبشی و پتانسیل}

انرژی جنبشی کل سیستم مجموع انرژی ارابه و آونگ است:
\[
    T = \frac{1}{2}M\dot{x}^2 + \frac{1}{2}m v_p^2
\]
\[
    T = \frac{1}{2}(M+m)\dot{x}^2 + \frac{1}{2}m \left( \dot{l}^2 + l^2\dot{\theta}^2 \right) + m\dot{x}\dot{l}\sin\theta + ml\dot{x}\dot{\theta}\cos\theta
\]
انرژی پتانسیل سیستم:
\[
    V = -mgl\cos\theta
\]

\subsection{تابع لاگرانژ}

تابع لاگرانژ سیستم ($\mathcal{L} = T - V$) تشکیل می‌شود:
\[
    \mathcal{L} = \frac{1}{2}(M+m)\dot{x}^2 + \frac{1}{2}m(\dot{l}^2 + l^2\dot{\theta}^2) + m\dot{x}\dot{l}\sin\theta + ml\dot{x}\dot{\theta}\cos\theta + mgl\cos\theta
\]

\subsection{نیروهای غیرمحافظه‌کار}
نیروهای تعمیم‌یافته:
\[
    Q_x = F_1 + F_2 - b_c\dot{x}
\]
\[
    Q_\theta = -b_p\dot{\theta}
\]

\subsection{اعمال معادلات اویلر-لاگرانژ}

\subsubsection{معادله حرکت ارابه ($x$)}
\[
    \frac{d}{dt}\left(\frac{\partial \mathcal{L}}{\partial \dot{x}}\right) - \frac{\partial \mathcal{L}}{\partial x} = Q_x
\]
با جایگذاری مشتقات و مرتب‌سازی ترم‌ها:
\[
    (M+m)\ddot{x} + ml\ddot{\theta}\cos\theta + m\ddot{l}\sin\theta + 2m\dot{l}\dot{\theta}\cos\theta - ml\dot{\theta}^2\sin\theta + b_c\dot{x} = F_1 + F_2
\]

\subsubsection{معادله حرکت آونگ ($\theta$)}
\[
    \frac{d}{dt}\left(\frac{\partial \mathcal{L}}{\partial \dot{\theta}}\right) - \frac{\partial \mathcal{L}}{\partial \theta} = Q_\theta
\]
پس از محاسبه مشتقات و حذف ترم‌های مشترک:
\[
    ml^2\ddot{\theta} + ml\ddot{x}\cos\theta + 2ml\dot{l}\dot{\theta} + mgl\sin\theta + b_p\dot{\theta} = 0
\]

\section{نیروهای تکیه‌گاهی ریل}

برای استخراج نیروهای تکیه‌گاهی $R_{x1}, R_{y1}, R_{y2}$، تعادل شبه‌استاتیکی کل مجموعه بررسی می‌شود:

\subsection{تعادل افقی}
نیروی خارجی مهارکننده ریل در راستای افق:
\[
    \sum F_x = 0 \implies R_{x1} - F_1 - F_2 + b_c\dot{x} = 0 \implies R_{x1} = F_1 + F_2 - b_c\dot{x}
\]

\subsection{تعادل قائم و گشتاور}
با فرض وزن ریل یکنواخت ($M_b g$) و نیروی عمودی ارابه $N = (M+m)g$ و همچنین $D$ طول ریل، از تعادل گشتاور حول تکیه‌گاه اول:
\[
    \sum M_1 = 0 \implies R_{y2} \cdot D - M_b g \frac{D}{2} - (M+m)g \cdot x = 0
\]
\[
    R_{y2} = \frac{M_b g}{2} + \frac{(M+m)g}{D} x
\]
تعادل نیروهای قائم:
\[
    \sum F_y = 0 \implies R_{y1} = \frac{M_b g}{2} + \frac{(M+m)g}{D} (D-x)
\]

\section{خطی‌سازی و نمایش فضای حالت}

حول نقطه کار $\theta \approx 0$ و فرض طول کابل اسمی ثابت $l(t) = l_0$ (در نتیجه $\dot{l} \approx 0, \ddot{l} \approx 0$ و $\sin\theta \approx \theta, \cos\theta \approx 1$)، معادلات خطی به فرم زیر در می‌آیند:
\[
    (M+m)\ddot{x} + ml_0\ddot{\theta} + b_c\dot{x} = F_1 + F_2
\]
\[
    ml_0^2\ddot{\theta} + ml_0\ddot{x} + b_p\dot{\theta} + mgl_0\theta = 0
\]

با استخراج شتاب‌ها از دستگاه فوق:
\[
    \ddot{x} = \frac{1}{M}\left(F_1 + F_2\right) - \frac{b_c}{M}\dot{x} + \frac{b_p}{Ml_0}\dot{\theta} + \frac{mg}{M}\theta
\]
\[
    \ddot{\theta} = -\frac{1}{Ml_0}\left(F_1 + F_2\right) + \frac{b_c}{Ml_0}\dot{x} - \left(\frac{b_p}{Ml_0^2} + \frac{b_p}{ml_0^2}\right)\dot{\theta} - \frac{(M+m)g}{Ml_0}\theta
\]

\subsection{فرم ماتریسی فضای حالت}

بردار حالت، ورودی و خروجی به صورت زیر تعریف می‌شوند:
\[
    \mathbf{x} = \begin{bmatrix} x \\ \dot{x} \\ \theta \\ \dot{\theta} \end{bmatrix}, \quad
    \mathbf{u} = \begin{bmatrix} F_1 \\ F_2 \end{bmatrix}, \quad
    \mathbf{y} = \begin{bmatrix} x \\ \theta \\ \Delta R_{x1} \\ \Delta R_{y1} \\ \Delta R_{y2} \end{bmatrix}
\]
خروجی نیروهای تکیه‌گاهی برحسب انحراف از مقدار تعادل بیان شده‌اند. ماتریس‌های فضای حالت ($\dot{\mathbf{x}} = A\mathbf{x} + B\mathbf{u}$ و $\mathbf{y} = C\mathbf{x} + D_m\mathbf{u}$):

\[
    A = \begin{bmatrix}
        0 & 1                 & 0                     & 0                                                       \\
        0 & -\dfrac{b_c}{M}   & \dfrac{mg}{M}         & \dfrac{b_p}{Ml_0}                                       \\
        0 & 0                 & 0                     & 1                                                       \\
        0 & \dfrac{b_c}{Ml_0} & -\dfrac{(M+m)g}{Ml_0} & -\left(\dfrac{b_p}{Ml_0^2} + \dfrac{b_p}{ml_0^2}\right)
    \end{bmatrix}
\]

\[
    B = \begin{bmatrix}
        0                & 0                \\
        \dfrac{1}{M}     & \dfrac{1}{M}     \\
        0                & 0                \\
        -\dfrac{1}{Ml_0} & -\dfrac{1}{Ml_0}
    \end{bmatrix}
\]

\[
    C = \begin{bmatrix}
        1                  & 0    & 0 & 0 \\
        0                  & 0    & 1 & 0 \\
        0                  & -b_c & 0 & 0 \\
        -\dfrac{(M+m)g}{D} & 0    & 0 & 0 \\
        \dfrac{(M+m)g}{D}  & 0    & 0 & 0
    \end{bmatrix}
\]

\[
    D_m = \begin{bmatrix}
        0 & 0 \\
        0 & 0 \\
        1 & 1 \\
        0 & 0 \\
        0 & 0
    \end{bmatrix}
\]

\subsection{تحلیل تکینگی کنترل‌پذیری در نقطه تعادل}

همان‌طور که در ماتریس‌های فضای حالت خطی‌سازی‌شده مشاهده می‌شود، متغیرهای مربوط به نرخ تغییر طول کابل (\lr{Reel Rate}) یعنی مشتقات \lr{$\dot{l}$} و \lr{$\ddot{l}$} به طور کامل از معادلات \lr{LTI} حذف شده‌اند و طول کابل تنها به صورت مقدار ثابت اسمی \lr{$l_0$} ظاهر شده است. دلیل این امر از منظر ریاضیاتِ خطی‌سازی (بسط تیلور مرتبه اول) این است که در نقطه کار قائم (\lr{$\theta \approx 0$})، ترم‌های کوپلینگ غیرخطی نظیر شتاب کوریولیس (\lr{$2m\dot{l}\dot{\theta}\cos\theta$}) و نیروی اینرسی ناشی از شتاب شعاعی (\lr{$m\ddot{l}\sin\theta$}) به دلیل ضرب شدن دو متغیر اغتشاشی در یکدیگر (جملات مرتبه دوم) و همچنین مقدار \lr{$\sin(0) = 0$}، به صفر میل می‌کنند. از دیدگاه فیزیکی، این پدیده بیانگر یک «تکینگی کنترل‌پذیری» (\lr{Controllability Singularity}) در نقطه تعادل آویزان است؛ بدین معنا که اگر آونگ کاملاً عمود باشد، بالا یا پایین کشیدن کابل هیچ‌گونه نیروی افقی یا گشتاوری برای نوسان ایجاد نمی‌کند. این نتیجه‌گیری اثبات می‌کند که کنترل‌کننده‌های خطی کلاسیک (نظیر \lr{LQR} یا \lr{Linear MPC}) ذاتاً قادر به بهره‌گیری از متغیر تغییر طول کابل جهت میرایی ارتعاشات حول نقطه تعادل نیستند. بنابراین، جهت استفاده‌ی مؤثر از قابلیت \lr{Reel Rate} به عنوان یک ابزار کنترلی (\lr{Active Damping})، عبور از مدل خطی و به‌کارگیری مدل کاملاً غیرخطی سیستم در قالب استراتژی‌های پیشرفته‌ای همچون کنترل‌کننده پیش‌بین مدل غیرخطی (\lr{NMPC}) امری الزامی و غیرقابل‌اجتناب است.
